\chapter{Literature Gaps}\label{ch:gaps}

The survey in Chapter~\ref{ch:survey} shows what zero-shot and supervised methods can do; this chapter
states what they leave open. We list five gaps that AgriAgent addresses, each with the papers it rests on
and the objective (Section~\ref{sec:objectives}) that responds to it, and then name the base paper that
our work builds on. Chapter~\ref{ch:proposed} then describes the system we built to close these gaps.

\begin{description}[style=nextline,leftmargin=0pt]
\item[G1. Supervised accuracy depends on labels from the same field.]
  Supervised networks separate crop and weed well~\cite{milioto2018realtime,lottes2018fcn,rahman2023cotton,dang2023yoloweeds},
  but they need large labelled, crop-specific datasets~\cite{hasan2021survey,lu2020datasets} and lose
  accuracy on new fields~\cite{weyler2023domain}. A method that needs \emph{no} labels from the target
  field has not been shown to be accurate enough for spraying. \emph{Response: O1, RQ1.}

\item[G2. Text-prompted zero-shot pipelines are untested on the spraying task.]
  Grounded SAM~\cite{ren2024groundedsam} is evaluated on general images, and the closest agricultural
  work~\cite{chong2025zeroshot} argues, without measuring it, that its text prompts are ill-defined for
  crop and weed; few-shot work~\cite{singh2025fewshot} reports that zero-shot prompts struggle with
  similar agricultural classes. No study we found measures Grounded SAM directly on crop-versus-weed
  photos with official labels. \emph{Response: O1, O3, RQ1.}

\item[G3. Segmentation metrics do not say what a sprayer needs to know.]
  Published zero-shot and supervised work reports pixel IoU or detection mAP~\cite{chong2025zeroshot,milioto2018realtime,rahman2023cotton}.
  A spot sprayer needs two different numbers: how much of the \emph{crop} lands on the spray map
  (damage) and how many weeds are hit (recall), together with the herbicide saved. None of the surveyed
  papers reports them for a zero-shot pipeline. \emph{Response: O3.}

\item[G4. No transparent safety layer between the model and the nozzle.]
  Smart sprayers act directly on the detector's output~\cite{partel2019sprayer}, and reviews identify
  the decision step as the weak point~\cite{slaughter2008review,gerhards2022sswm}. Vision-language
  models are useful second opinions but unreliable alone~\cite{nasir2026vlmweed}, and LLM agents act on
  their own conclusions~\cite{wang2024agents}. No surveyed system places a deterministic, testable
  gate, which withholds anything uncertain, between the models and the spraying decision.
  Row-based methods~\cite{bah2018unsupervised,louargant2018rowcrops} show that simple agronomic
  structure can replace labels, but they have not been combined with a zero-shot detector as a safety
  layer. \emph{Response: O2, O4, RQ2, RQ3.}

\item[G5. Advice from language models is not grounded in local regulation.]
  Reviews of LLMs in agriculture~\cite{kuska2024llm,li2025llmagri,yin2025foundation} discuss their
  promise for advisory work but also their reliability problems. A herbicide recommendation must come
  from an approved national source with legal doses and waiting periods, here the ICAR
  bulletin~\cite{choudhury2016icar}, not from the model. \emph{Response: O5, O6.}
\end{description}

\section{Base paper}
The pipeline AgriAgent implements is Grounded SAM~\cite{ren2024groundedsam}, with SAM~2~\cite{ravi2025sam2}
in place of SAM. Grounded SAM is an arXiv technical report, so we take Chong et
al.~\cite{chong2025zeroshot} (IEEE/RSJ IROS 2025) as the peer-reviewed reference point: it addresses
the same task, zero-shot, and names our pipeline's weakness. Our work fills gaps G2--G4 for that
pipeline: we measure its weakness on the sprayer's metrics and test how much a transparent layer of
rules and a gate recovers.
